\bookchapter{第二十三章：编程思维与自动化启蒙}

\chapterintro{前两章，咱们给手机筑起了安全的“边界墙”，也见识了它在墙内日益聪明的“AI大脑”。现在，让我们解锁手机的终极潜力之一：让它像程序员一样“思考”，并自动完成你设定的连续任务。你可能会说：“我又不会写代码，这跟我有什么关系？” 放心，本章讲的根本不需要你写一行代码——现代手机已经把“编程”的门槛降到了拖拖拽拽就能搞定的程度。}

\booksection{一、自动化到底是个啥？}

用大白话说，自动化就是：\textbf{让手机在某个条件满足时，自动帮你做一连串事。}你不需要每次都手动操作，也不用记住该做什么——手机会自己判断时机，自己执行。

举个例子：每天早上你醒来，要手动关掉勿扰模式、打开微信看看消息、打开天气App看今天穿什么、打开网易云放首歌。如果把这些步骤打包成一个自动化流程，设置成“每天早上7点自动执行”，那你醒了只需要拿起手机，一切都已经准备好了。是不是想想就很爽？

自动化本质上就是“如果……那么……”的逻辑链条，编程里的if-else。学明白了这个思维方式，哪怕你一辈子不写代码，处理问题的条理性也会强很多。核心要素：

\begin{enumerate}
  \item \textbf{触发条件}：什么时候开始？可以是"每天早上7点"、"当我连接到家里的Wi-Fi"、"当我到达某个地点"、"当手机电量低于20\%"。
  \item \textbf{执行动作}：具体做什么？可以是发一条微信、打开某个App、调整音量、切换壁纸、播放音乐、发送通知……几乎所有你能在手机上手动做的事，自动化都能做。
  \item \textbf{中间逻辑}（可选但很强大）：如果情况A就做X，否则做Y。比如"如果是工作日就播放新闻播客，周末就播放音乐"，这叫条件分支。
\end{enumerate}

\booksection{二、iOS“快捷指令”——苹果用户的自动化神器}

苹果的“快捷指令”（Shortcuts）App是自带的应用，不需要下载。它的核心概念就是把各种操作像搭积木一样拼起来。先说几个关键概念：

\begin{itemize}
  \item \textbf{"快捷指令"}：一个完整的自动化流程，可以手动点击运行，也可以通过Siri语音唤醒。比如你说"嘿Siri，回家模式"，它就自动执行你设定好的一连串操作。
  \item \textbf{"自动化"}：由触发条件自动运行的指令，不需要你手动触发。比如"当手机连接到车载蓝牙时，自动打开地图App"。自动化在"自动化"标签页下设置，和"快捷指令"是分开的。
  \item \textbf{"动作"}：最小的执行单元。打开App是一个动作，发送消息是一个动作，获取天气是一个动作，调整亮度是一个动作。你可以把成百上千个动作串在一起。
\end{itemize}

给你看几个实实在在好用的例子，你跟着做一遍就明白了：

\booksubsection{实例一：一键“早安流程”}

每天早上点一下（或对Siri说“早上好”），手机会自动：获取今天的天气朗读天气和气温显示今天的日历安排打开你常用的音乐App开始播放。操作步骤：新建快捷指令添加动作“获取当前天气”添加“朗读文本”添加“获取日历事件”添加“显示提醒”添加“打开App”选网易云音乐。就这么简单，全程拖拽。

\booksubsection{实例二：自动“省电模式”}

在“自动化”标签页，新建自动化选“电池电量”设“低于30\%”添加动作“设定低电量模式为开启”、“降低屏幕亮度到30\%”、“关闭蓝牙”。从此手机电量低于30\%就自动进入省电状态，不用你操心。

\booksubsection{实例三：扫码自动记账}

每次你用微信/支付宝扫码支付后，手机会读取支付通知，自动把金额和商家记录到一个备忘录或Numbers表格里。这需要用到快捷指令的“自动化”里的“收到短信/通知”触发条件（iOS 17+支持读取钱包交易），配合“追加到Numbers”动作。对于想记账又懒得手动记的人来说，简直是救星。

还有一个很多人不知道的巧用：你可以把快捷指令添加到主屏幕，变成一个App图标。比如做一个“打车回家”的快捷指令，图标放桌面上，点一下就直接打开滴滴并把目的地设成家，比你自己打开App慢慢输入快多了。

\booksection{三、安卓自动化——Tasker和它的朋友们}

安卓这边，情况稍微复杂一点，但天花板也更高。先说简单的：

\booksubsection{入门级：自动点击器}

很多安卓手机自带“自动点击器”或者可以在应用商店下载。它的原理超简单：你手动操作一遍，它把你点击的坐标和时间间隔录下来，然后可以高速回放。比如某个App每天要领签到奖励，你录一遍点击流程，以后每天自动跑一遍就完事了。缺点是：它只认坐标，如果界面布局变了（比如弹了个广告），它就点歪了。

\booksubsection{进阶级：Tasker——安卓自动化之王}

Tasker是安卓平台上公认最强的自动化工具，常年霸榜。说实话它的学习曲线有点陡，界面不算好看，但一旦上手，你会觉得手机突然变得无所不能。Tasker的核心也是“触发任务”的模式：

\begin{itemize}
  \item \textbf{配置文件（Profile）}：定义"什么时候触发"。触发条件可以是时间、地点、Wi-Fi连接状态、蓝牙连接、收到特定通知、电池状态、屏幕朝向……Tasker支持的触发条件多到离谱，有几十种。
  \item \textbf{任务（Task）}：定义"触发后做什么"。一个任务可以包含多个动作，按顺序执行。动作类型包括：调整系统设置（音量、亮度、Wi-Fi、蓝牙、飞行模式）、发送短信、HTTP请求、执行Shell命令、控制媒体播放、弹出对话框等等。
\end{itemize}

给你几个Tasker的经典实用场景：

\begin{itemize}
  \item \textbf{到家自动开Wi-Fi}：配置触发条件为"连接到特定基站"（因为你在家附近会连到固定的基站）任务"开启Wi-Fi"。不用GPS也不耗电，手机一进小区范围就自动开Wi-Fi。
  \item \textbf{会议模式}：触发条件为"日历事件标题包含'会议'"任务"设置静音模式、关闭媒体音量、自动回复微信（需要插件配合）"。再也不用担心开会时手机突然响了。
  \item \textbf{深夜护眼}：触发条件为"时间22:00-06:00"任务"开启护眼模式、降低亮度到20\%、将屏幕设为黑白（灰度模式）"。黑白屏幕对刷短视频的欲望有奇效——画面一灰，你就不想刷了。
  \item \textbf{耳机插入自动播客}：触发条件为"有线耳机插入/蓝牙耳机连接"任务"打开小宇宙App并开始播放最新一集"。戴上耳机就能听，不用解锁手机找App。
\end{itemize}

如果你觉得Tasker上手太难，可以试试它的简化版兄弟：\textbf{MacroDroid}。界面友好很多，模板丰富，适合不想折腾的用户。另外现在很多国产安卓系统（小米的MIUI、华为的鸿蒙、OPPO的ColorOS）都自带了“场景”或“智能助手”功能，虽然没有Tasker那么强大，但日常自动化完全够用。

\booksection{四、自动化思维的日常应用}

学会自动化最大的收获，其实不是省那点时间，而是\textbf{改变你思考问题的方式}。一旦你习惯了“如果发生X，就自动做Y”的思维模式，你会发现生活中一大堆重复劳动都可以被干掉。比如：

\begin{itemize}
  \item 每次拍完照片，自动备份到网盘然后把本地照片压缩成小尺寸节省空间。
  \item 每天早上自动从日历读取今天的安排，生成一段文字通过微信发给家人。
  \item 收到快递签收短信后，自动在备忘录里添加一条待办"取快递"。
  \item 手机电量充满时自动发出语音提醒拔掉充电器（避免过充伤电池）。
  \item 每天下班时间到了，如果还在公司附近（通过基站定位），自动提醒"该下班了"。
\end{itemize}

这些听起来很高级，但其实你只要掌握了“触发+动作”这两个基本概念，配合快捷指令或Tasker，都能做出来。关键是\textbf{你要先意识到“这件事可以自动化”}——很多人不是不会，是根本没往那方面想。

\booksection{五、踩坑指南：自动化的常见问题}

搞自动化也不是一帆风顺的，有几个坑我帮你先踩过了：

\begin{itemize}
  \item \textbf{位置触发可能不灵}：iOS的"到达/离开位置"触发用的是地理围栏技术，准确度取决于GPS信号和系统优化。有时候你都到家十分钟了才触发。建议用Wi-Fi连接或基站触发代替位置触发，更可靠也更省电。
  \item \textbf{快捷指令的"自动化"需要授权}：iOS里，"自动化"标签下的触发大多数需要你手动确认才能运行（iOS的安全机制）。但Wi-Fi连接、蓝牙连接、闹钟停止这几个触发是可以真正全自动运行的，不需要确认。
  \item \textbf{Tasker需要开启无障碍权限}：很多高级功能（比如模拟点击）需要无障碍服务权限，这会让手机在设置里显示一个常驻通知。这是安卓的安全设计，没法去掉。
  \item \textbf{别把自动化搞得太复杂}：刚开始做自动化，从"2-3个动作"的小流程开始。别一上来就搞一个20步的超级流程——调试起来会疯掉的。
\end{itemize}

\textbf{小贴士：自动化入门的正确姿势}\par

先花5分钟想一下：你每天在手机上重复最多的3件事是什么？比如每天早上的例行操作、每次连蓝牙耳机后的习惯、每次出门前的设置……从这些高频小事入手，各做一个简单的自动化。哪怕每个自动化只省你5秒钟，日积月累也是可观的。而且做出来之后的满足感，真的会上瘾。

\begin{tcbpractice}{练一练：三个入门自动化挑战}
选一个最适合你的完成，做完你就入门了：

\begin{enumerate}
  \item \textbf{挑战一：早安自动化}：iOS：在快捷指令App里新建快捷指令添加"获取当前天气""朗读文本""打开App选网易云"。完成后你可以对Siri说"早上好"来触发。安卓：在Tasker里新建任务添加"获取天气""朗读""打开App"设置触发时间为每天7:00。
  \item \textbf{挑战二：低电量自动省电}：iOS：自动化标签新建"电池电量"低于25\%添加"设定低电量模式开"+"设定亮度25\%"。安卓：Tasker新建Profile状态电源电量0-25\%任务"省电模式开"+"亮度调低"。
  \item \textbf{挑战三：连接车载蓝牙自动导航}：iOS：自动化新建"蓝牙"选择你的车载蓝牙添加"打开App选高德地图"+"媒体音量调至80\%"。安卓：TaskerProfile状态蓝牙已连接选择车载蓝牙任务"打开高德"+"媒体音量调高"。
\end{enumerate}

\end{tcbpractice}

\begin{tcbinfo}{案例 3：自动化"早安流程"}
\begin{itemize}
  \item \textbf{场景}：每天 7:00 自动播报天气 + 关闭静音 + 打开新闻。
  \item \textbf{步骤}：iOS: 快捷指令自动化时间触发添加"获取天气"+"设定音量"+"打开 App";Android: MacroDroid触发器(时间 7:00)动作(媒体音量/打开 URL/TTS 朗读天气)。
  \item \textbf{预期结果}：免安装即可实现，降低上手门槛。
\end{itemize}

\end{tcbinfo}
